\documentclass[11pt]{article}

\usepackage[margin=1.1in]{geometry}
\usepackage{amsmath,amssymb,amsthm,mathtools}
\usepackage{microtype}
\usepackage{booktabs}
\usepackage{enumitem}
\usepackage{hyperref}

\hypersetup{
    colorlinks=true,
    linkcolor=black,
    citecolor=black,
    urlcolor=black
}

\newtheorem{theorem}{Theorem}[section]
\newtheorem{corollary}[theorem]{Corollary}
\newtheorem{proposition}[theorem]{Proposition}
\newtheorem{lemma}[theorem]{Lemma}
\newtheorem{remark}[theorem]{Remark}

\newcommand{\R}{\mathbb{R}}
\newcommand{\C}{\mathbb{C}}
\newcommand{\LiSus}{\chi^{\mathrm{Li}}}
\newcommand{\Rint}{R_{\mathrm{int}}}
\newcommand{\Ebar}{\overline E}

\title{
Matched-Scale Li Susceptibility and Centered-Primitive Transfer
of Zero-Counting Discrepancy
}

\author{Michael Elsbernd}

\date{Research draft}

\begin{document}

\maketitle

\begin{abstract}
We study a horizontal zero-coordinate response associated with the
Li kernel
\[
f_n(z)
=
1-
\left(
1-\frac1z
\right)^n.
\]
The corresponding quadratic susceptibility is
\[
\LiSus_n(\gamma)
=
2\operatorname{Re}
f_n''
\left(
\frac12+i\gamma
\right).
\]
In the matched-scale regime \(n/H\asymp1\), we derive compact-uniform
asymptotic profiles for the susceptibility and its first two
\(\gamma\)-derivatives.

We then consider the global zero-counting discrepancy
\[
E(t)=N(t)-M(t)
\]
on a fixed relative window
\[
[H,(1+\eta)H].
\]
After subtracting the window mean of \(E\), we introduce its centered
primitive \(F\). An exact second-order integration-by-parts identity
transfers the response
\[
\Rint(n;H,\eta)
=
-
\int_H^{(1+\eta)H}
E(t)\LiSus_n'(t)\,dt
\]
from the rough counting discrepancy to \(F\) and the second derivative
of the Li susceptibility.

Using the classical unconditional bound
\[
S_1(T)=O(\log T)
\]
for the first antiderivative of the argument term in the
Riemann--von Mangoldt formula, we prove
\[
\Rint(n;H,\eta)
=
O_{K,\eta}
\left(
\frac{\log H}{H^3}
\right)
\]
uniformly for \(n/H\) in a fixed compact subset
\(K\Subset(0,\infty)\). Consequently,
\[
H^2\Rint(n;H,\eta)\to0.
\]
Under the Riemann Hypothesis, the classical sharper bound for
\(S_1\) gives the logarithmic refinement
\[
\Rint(n;H,\eta)
=
O_{K,\eta}
\left(
\frac{\log H}
{H^3(\log\log H)^2}
\right).
\]

The result is a response-asymptotic statement and is compatible with
failure of the Riemann Hypothesis. Numerical zero-window experiments
are presented only as finite-height motivation for the centered-
primitive mechanism.
\end{abstract}

\section{Introduction}

Let
\[
f_n(z)
=
1-
\left(
1-\frac1z
\right)^n,
\qquad n\ge1.
\]
This is the individual-zero kernel appearing in Li-type zero sums.
The relation between Li's coefficients and the Guinand--Weil explicit
formula is classical; see Bombieri and Lagarias
\cite{BombieriLagarias1999}. Our purpose is not to formulate another
Li criterion or another Weil positivity criterion.

Instead, we isolate a local horizontal response of \(f_n\) at the
critical line and study how the global zero-counting discrepancy acts
through the resulting response kernel.

For an analytic function \(f\), define the symmetric horizontal orbit
response
\[
R_f(\delta,\gamma)
=
\sum_{\epsilon=\pm1}
\sum_{\sigma=\pm1}
f
\left(
\frac12+\epsilon\delta+i\sigma\gamma
\right)
-
2
\sum_{\sigma=\pm1}
f
\left(
\frac12+i\sigma\gamma
\right).
\]
Taylor expansion in the horizontal displacement gives
\[
R_f(\delta,\gamma)
=
\delta^2\chi_f(\gamma)
+
O(\delta^4),
\]
where
\[
\chi_f(\gamma)
=
2\operatorname{Re}
f''
\left(
\frac12+i\gamma
\right).
\]
For the Li kernel, we write
\[
\LiSus_n(\gamma)
=
2\operatorname{Re}
f_n''
\left(
\frac12+i\gamma
\right).
\]

Our main object is the ordinary-integral response
\[
\Rint(n;H,\eta)
=
-
\int_H^{(1+\eta)H}
E(t)\LiSus_n'(t)\,dt,
\]
where
\[
E(t)=N(t)-M(t)
\]
is the discrepancy between the classical nontrivial-zero counting
function and the leading Riemann--von Mangoldt model.

The central theorem is the following.

\begin{theorem}[Unconditional global Li-response suppression]
\label{thm:main}
Fix \(\eta>0\) and a compact interval
\[
K\Subset(0,\infty).
\]
Uniformly over positive integers \(n\) satisfying
\[
\frac nH\in K,
\]
one has
\[
\Rint(n;H,\eta)
=
O_{K,\eta}
\left(
\frac{\log H}{H^3}
\right)
\]
as \(H\to\infty\). Consequently,
\[
H^2\Rint(n;H,\eta)\to0
\]
uniformly on compact positive matched-scale ranges.
\end{theorem}

The proof is short once the matched-scale derivative hierarchy of
the Li susceptibility has been established. It combines four ingredients.

First, the classical Riemann--von Mangoldt formula gives
\[
E(t)=S(t)+O(1/t).
\]
Second, the classical unconditional estimate
\[
S_1(T)=O(\log T),
\qquad
S_1(T)=C+\int_0^T S(t)\,dt,
\]
controls an antiderivative of the counting discrepancy.
Third, the Li susceptibility satisfies a matched-scale second-derivative
bound
\[
\int_H^{(1+\eta)H}
|\LiSus_n''(t)|\,dt
=
O_{K,\eta}(H^{-3}).
\]
Finally, centering \(E\) on the window produces a primitive with
vanishing boundary values, allowing a second integration by parts.

Under the Riemann Hypothesis, the sharper classical estimate
\[
S_1(T)
=
O
\left(
\frac{\log T}
{(\log\log T)^2}
\right)
\]
immediately yields the following refinement.

\begin{corollary}[RH-conditional logarithmic sharpening]
\label{cor:rh}
Assume the Riemann Hypothesis. Under the hypotheses of
Theorem~\ref{thm:main},
\[
\Rint(n;H,\eta)
=
O_{K,\eta}
\left(
\frac{\log H}
{H^3(\log\log H)^2}
\right).
\]
\end{corollary}

The unconditional theorem is compatible with zeros away from the
critical line. Indeed, the scalar counting function \(N(T)\) records
ordinates and multiplicities but discards the horizontal coordinates
of zeros. The theorem should therefore be interpreted as a
zero-counting response estimate, not as an RH criterion.

\section{Classical zero-counting input}

Let \(N(T)\) count the nontrivial zeros
\[
\rho=\beta+i\gamma
\]
with
\[
0<\gamma\le T,
\]
counted with multiplicity.

Set
\[
M(T)
=
\frac{T}{2\pi}
\log\frac{T}{2\pi}
-
\frac{T}{2\pi}
+
\frac78.
\]
The classical Riemann--von Mangoldt formula gives
\[
N(T)
=
M(T)+S(T)+R(T),
\]
where
\[
R(T)=O(1/T).
\]
Thus
\[
E(T)
:=
N(T)-M(T)
=
S(T)+R(T).
\]

We use ordinary integrals of \(E\). Values chosen for \(S(T)\) at
individual zero ordinates therefore do not affect the main argument.
Any change of \(S(t)\) on the discrete zero set changes neither the
ordinary integrals used below nor the antiderivative \(S_1\) up to the
same harmless convention.

Define
\[
S_1(T)
=
C+\int_0^T S(t)\,dt.
\]
We use the classical unconditional estimate
\begin{equation}
\label{eq:s1-unconditional}
S_1(T)=O(\log T).
\end{equation}
Under RH, we use
\begin{equation}
\label{eq:s1-rh}
S_1(T)
=
O
\left(
\frac{\log T}
{(\log\log T)^2}
\right).
\end{equation}
These estimates belong to the classical theory of the argument of the
Riemann zeta function; see, for example, the literature summarized in
Wakasa \cite{WakasaS1}.

\section{The Li susceptibility}

\subsection{Exact formula}

For
\[
f_n(z)
=
1-
\left(
1-\frac1z
\right)^n,
\]
define
\[
\LiSus_n(\gamma)
=
2\operatorname{Re}
f_n''
\left(
\frac12+i\gamma
\right).
\]

Set
\[
\theta_\gamma
=
2\arctan
\left(
\frac1{2\gamma}
\right).
\]

A direct calculation gives
\begin{equation}
\label{eq:chi-exact}
\LiSus_n(\gamma)
=
-
\frac{2n}
{
(\gamma^2+1/4)^2
}
\left[
n\cos(n\theta_\gamma)
+
2\gamma\sin(n\theta_\gamma)
\right].
\end{equation}

\subsection{Leading matched-scale profile}

Let
\[
q=\frac n\gamma.
\]
Define
\begin{equation}
\label{eq:Psi}
\Psi_{\mathrm{Li}}(q)
=
-2q
\left(
q\cos q+2\sin q
\right).
\end{equation}

\begin{proposition}
\label{prop:chi-profile}
For every compact set
\[
K\Subset(0,\infty),
\]
uniformly for
\[
q=\frac n\gamma\in K,
\]
one has
\[
\gamma^2\LiSus_n(\gamma)
=
\Psi_{\mathrm{Li}}(q)
+
O_K(\gamma^{-2}).
\]
\end{proposition}

\begin{proof}
Write
\[
x=\gamma^{-1}
\]
and
\[
d(x)=1+\frac{x^2}{4}.
\]
Then
\[
n\theta_\gamma
=
\frac qx
2\arctan\frac x2.
\]
The function
\[
\frac2x\arctan\frac x2
\]
is analytic in \(x^2\) near \(x=0\) and equals \(1\) at the origin.
Substitution into the exact expression
\eqref{eq:chi-exact} produces an analytic expansion in \(q\) and \(x^2\)
on a neighborhood of every compact positive \(q\)-set. Setting
\(x=0\) yields \(\Psi_{\mathrm{Li}}(q)\), and the first omitted power is
\(x^2\). This proves the compact-uniform estimate.
\end{proof}

\subsection{First derivative profile}

Define
\begin{equation}
\label{eq:J}
J_{\mathrm{Li}}(q)
=
2q^3\sin q
-
12q^2\cos q
-
12q\sin q.
\end{equation}

The first derivative must be obtained from the exact scaled analytic
representation rather than by formally differentiating the error term
in Proposition~\ref{prop:chi-profile}.

Let
\[
x=\frac1\gamma,
\qquad
q=\frac n\gamma,
\qquad
d(x)=1+\frac{x^2}{4},
\]
and
\[
\phi(q,x)
=
\frac qx
2\arctan\frac x2.
\]
Direct differentiation of the exact susceptibility gives
\[
\gamma^3\LiSus_n'(\gamma)
=
G(q,x),
\]
where
\[
\begin{aligned}
G(q,x)
={}&
\frac{8q}{d^3}
\left(
q\cos\phi+2\sin\phi
\right)
\\
&-
\frac{2q}{d^2}
\left[
\left(
\frac{q^2}{d}+2
\right)
\sin\phi
-
\frac{2q}{d}\cos\phi
\right].
\end{aligned}
\]
The function \(G\) is analytic in \(q\) and \(x^2\) near compact
positive \(q\)-sets, and
\[
G(q,0)=-J_{\mathrm{Li}}(q).
\]

Therefore:

\begin{proposition}
\label{prop:first-derivative}
Uniformly for
\[
q=\frac n\gamma
\]
in compact positive sets,
\[
\gamma^3\LiSus_n'(\gamma)
=
-J_{\mathrm{Li}}(q)
+
O_K(\gamma^{-2}).
\]
\end{proposition}

\subsection{Second derivative profile}

Define
\begin{equation}
\label{eq:K2-definition}
K_{\mathrm{Li}}^{(2)}(q)
=
3J_{\mathrm{Li}}(q)
+
qJ_{\mathrm{Li}}'(q).
\end{equation}

From \eqref{eq:J},
\[
J_{\mathrm{Li}}'(q)
=
2q^3\cos q
+
18q^2\sin q
-
36q\cos q
-
12\sin q.
\]
Hence
\begin{equation}
\label{eq:K2}
\boxed{
K_{\mathrm{Li}}^{(2)}(q)
=
2q^4\cos q
+
24q^3\sin q
-
72q^2\cos q
-
48q\sin q.
}
\end{equation}

We again use the exact scaled analytic representation.

Since
\[
\LiSus_n'(\gamma)
=
\gamma^{-3}G(q,x),
\]
where
\[
q=\frac n\gamma,
\qquad
x=\frac1\gamma,
\]
and
\[
\frac{dq}{d\gamma}
=
-\frac q\gamma,
\qquad
\frac{dx}{d\gamma}
=
-\frac1{\gamma^2},
\]
we obtain
\[
\gamma^4\LiSus_n''(\gamma)
=
-3G(q,x)
-
qG_q(q,x)
-
xG_x(q,x).
\]
At \(x=0\),
\[
G(q,0)=-J_{\mathrm{Li}}(q),
\]
\[
G_q(q,0)=-J_{\mathrm{Li}}'(q),
\]
and analyticity in \(x^2\) gives
\[
xG_x(q,x)=O_K(x^2).
\]

Thus:

\begin{proposition}
\label{prop:second-derivative}
Uniformly for
\[
q=\frac n\gamma
\]
in compact positive sets,
\[
\boxed{
\gamma^4\LiSus_n''(\gamma)
=
K_{\mathrm{Li}}^{(2)}(q)
+
O_K(\gamma^{-2}).
}
\]
\end{proposition}

\section{Relative-window kernel bounds}

Fix
\[
\eta>0
\]
and
\[
K=[c_-,c_+]
\Subset(0,\infty).
\]
Suppose
\[
\frac nH\in K.
\]
Set
\[
B=(1+\eta)H.
\]

For
\[
t=Hu,
\qquad
1\le u\le1+\eta,
\]
we have
\[
\frac nt
=
\frac{n/H}{u}
\in
K_\eta
:=
\left[
\frac{c_-}{1+\eta},
c_+
\right].
\]
The set \(K_\eta\) is compact and contained in \((0,\infty)\).

Proposition~\ref{prop:chi-profile} therefore implies
\[
\LiSus_n(t)
=
O_{K,\eta}(H^{-2})
\]
uniformly on the window. In particular,
\begin{equation}
\label{eq:endpoint-scale}
|\LiSus_n(B)-\LiSus_n(H)|
=
O_{K,\eta}(H^{-2}).
\end{equation}

Similarly, Proposition~\ref{prop:second-derivative} gives
\[
\LiSus_n''(t)
=
O_{K,\eta}(H^{-4})
\]
uniformly. Since
\[
B-H=\eta H,
\]
we obtain
\begin{equation}
\label{eq:second-variation-scale}
\int_H^B
|\LiSus_n''(t)|\,dt
=
O_{K,\eta}(H^{-3}).
\end{equation}

More precisely, if
\[
\frac nH\to c>0,
\]
then
\[
H^3
\int_H^{(1+\eta)H}
|\LiSus_n''(t)|\,dt
\to
\int_1^{1+\eta}
u^{-4}
\left|
K_{\mathrm{Li}}^{(2)}
\left(
\frac cu
\right)
\right|\,du.
\]

\section{Centered discrepancy primitive}

Define the window mean
\[
\Ebar
=
\frac1{B-H}
\int_H^B
E(s)\,ds.
\]

Using
\[
E=S+R,
\]
we have
\[
\int_H^B E(s)\,ds
=
S_1(B)-S_1(H)
+
\int_H^B R(s)\,ds.
\]
By \eqref{eq:s1-unconditional},
\[
S_1(B)-S_1(H)
=
O_\eta(\log H).
\]
Also,
\[
\int_H^B|R(s)|\,ds
=
O_\eta(1).
\]
Since
\[
B-H=\eta H,
\]
we obtain
\begin{equation}
\label{eq:Ebar}
\boxed{
\Ebar
=
O_\eta
\left(
\frac{\log H}{H}
\right).
}
\end{equation}

Now define
\[
F(t)
=
\int_H^t
\left(
E(s)-\Ebar
\right)\,ds.
\]
Then
\[
F(H)=F(B)=0.
\]

Let
\[
A(t)
=
\int_{t_0}^t E(s)\,ds
\]
for fixed \(t_0>0\). Directly,
\[
\boxed{
F(t)
=
A(t)-A(H)
-
\frac{t-H}{B-H}
\left(
A(B)-A(H)
\right).
}
\]
Thus \(F\) is the secant-detrended increment of an antiderivative of
the counting discrepancy.

Since
\[
A(t)
=
S_1(t)
+
\int_{t_0}^t R(s)\,ds
+
C_0,
\]
the constant \(C_0\) cancels. On the relative window,
\[
\sup_{H\le t\le B}
|S_1(t)|
=
O_\eta(\log H),
\]
and all integrated remainder increments are \(O_\eta(1)\). Hence
\begin{equation}
\label{eq:Fbound}
\boxed{
\|F\|_{\infty,[H,B]}
=
O_\eta(\log H).
}
\end{equation}

\section{Exact response transfer}

Define
\[
\Rint(n;H,\eta)
=
-
\int_H^B
E(t)\LiSus_n'(t)\,dt.
\]

Since
\[
F'(t)
=
E(t)-\Ebar
\]
almost everywhere,
\[
E(t)=\Ebar+F'(t)
\]
almost everywhere. Therefore
\[
\begin{aligned}
\Rint
&=
-\Ebar
\int_H^B
\LiSus_n'(t)\,dt
-
\int_H^B
F'(t)\LiSus_n'(t)\,dt.
\end{aligned}
\]
Because
\[
F(H)=F(B)=0,
\]
and because \(F\) is absolutely continuous while \(\LiSus_n'\) is
\(C^1\) on the window, the ordinary absolutely-continuous
integration-by-parts formula yields
\begin{equation}
\label{eq:transfer}
\boxed{
\Rint
=
-\Ebar
\left(
\LiSus_n(B)-\LiSus_n(H)
\right)
+
\int_H^B
F(t)\LiSus_n''(t)\,dt.
}
\end{equation}

We call the first term
\[
R_{\mathrm{mean}}
=
-\Ebar
\left(
\LiSus_n(B)-\LiSus_n(H)
\right)
\]
and the second
\[
R_{\mathrm{prim}}
=
\int_H^B
F(t)\LiSus_n''(t)\,dt.
\]

Thus
\[
\Rint
=
R_{\mathrm{mean}}
+
R_{\mathrm{prim}}.
\]

\section{Proof of the main theorem}

We now prove Theorem~\ref{thm:main}.

From \eqref{eq:Ebar} and \eqref{eq:endpoint-scale},
\[
\begin{aligned}
|R_{\mathrm{mean}}|
&\le
|\Ebar|
\,
|\LiSus_n(B)-\LiSus_n(H)|
\\
&=
O_\eta
\left(
\frac{\log H}{H}
\right)
O_{K,\eta}(H^{-2}).
\end{aligned}
\]
Therefore
\begin{equation}
\label{eq:mean-bound}
R_{\mathrm{mean}}
=
O_{K,\eta}
\left(
\frac{\log H}{H^3}
\right).
\end{equation}

For the primitive channel, \eqref{eq:Fbound} and
\eqref{eq:second-variation-scale} give
\[
\begin{aligned}
|R_{\mathrm{prim}}|
&\le
\|F\|_\infty
\int_H^B
|\LiSus_n''(t)|\,dt
\\
&=
O_\eta(\log H)
O_{K,\eta}(H^{-3}).
\end{aligned}
\]
Hence
\begin{equation}
\label{eq:primitive-bound}
R_{\mathrm{prim}}
=
O_{K,\eta}
\left(
\frac{\log H}{H^3}
\right).
\end{equation}

Combining \eqref{eq:transfer}, \eqref{eq:mean-bound}, and
\eqref{eq:primitive-bound}, we obtain
\[
\boxed{
\Rint(n;H,\eta)
=
O_{K,\eta}
\left(
\frac{\log H}{H^3}
\right).
}
\]
Multiplying by \(H^2\) gives
\[
H^2\Rint(n;H,\eta)
=
O_{K,\eta}
\left(
\frac{\log H}{H}
\right)
\to0.
\]
This proves Theorem~\ref{thm:main}.

No assumption of the Riemann Hypothesis is used.

\section{The RH-conditional refinement}

Assume RH.

Replacing \eqref{eq:s1-unconditional} by
\eqref{eq:s1-rh} gives
\[
\Ebar
=
O_\eta
\left(
\frac{\log H}
{H(\log\log H)^2}
\right)
\]
and
\[
\|F\|_\infty
=
O_\eta
\left(
\frac{\log H}
{(\log\log H)^2}
\right).
\]
The susceptibility estimates are unchanged. Substitution into
\eqref{eq:transfer} yields
\[
\boxed{
\Rint(n;H,\eta)
=
O_{K,\eta}
\left(
\frac{\log H}
{H^3(\log\log H)^2}
\right).
}
\]
This proves Corollary~\ref{cor:rh}.

\section{Normalized response shares}

The main theorem may be converted into a normalized-share statement
only when a competing response channel has a lower bound at the
natural susceptibility scale.

\begin{proposition}
\label{prop:abstract-share}
Suppose along a matched-scale sequence that a competing response
channel \(R_{\mathrm{comp}}\) satisfies
\[
|R_{\mathrm{comp}}|
\ge
cH^{-2}
\]
for some fixed \(c>0\). Then
\[
\frac{|\Rint|}
{
|R_{\mathrm{comp}}|+|\Rint|
}
=
O_{K,\eta,c}
\left(
\frac{\log H}{H}
\right)
\to0.
\]
\end{proposition}

\begin{proof}
By Theorem~\ref{thm:main},
\[
|\Rint|
\le
C_{K,\eta}
\frac{\log H}{H^3}
\]
for sufficiently large \(H\). Therefore
\[
\frac{|\Rint|}
{
|R_{\mathrm{comp}}|+|\Rint|
}
\le
\frac{
C_{K,\eta}\log H/H^3
}{
cH^{-2}
}
=
\frac{C_{K,\eta}}{c}
\frac{\log H}{H}.
\]
\end{proof}

For a Stieltjes decomposition, the atom channel at the left endpoint is
\[
R_{\mathrm{atom}}
=
\LiSus_n(H)\Delta N(H).
\]
Kernel nondegeneracy alone does not imply atom existence. We therefore
state the atom specialization with an explicit jump hypothesis.

\begin{corollary}
\label{cor:atom-share}
Suppose
\[
\frac nH\in K,
\qquad
|\Psi_{\mathrm{Li}}(n/H)|
\ge\delta>0,
\]
and
\[
\Delta N(H)=m_H\ge m_0>0.
\]
Then
\[
|R_{\mathrm{atom}}|
\gg_{K,\delta,m_0}
H^{-2},
\]
and the normalized interior response share tends to zero.
\end{corollary}

\section{Finite-height numerical motivation}

The centered-primitive formulation was motivated by finite-height
experiments using computed critical-line zero ordinates.

A first block used the first \(500\) computed ordinates. A disjoint
higher-index block used ordinates with indices \(501\) through \(1500\).
The second block was evaluated against a component-share principal
axis frozen from the first block before the higher-block response
shares were generated.

The numerical response decomposition exhibited a small interior
absolute share and a component-share cloud strongly concentrated near
the atom--endpoint exchange direction. A centered-primitive transfer
certificate was then evaluated numerically and displayed substantial
finite-height compression relative to the first absolute transfer
scale.

These computations are not proof dependencies for
Theorem~\ref{thm:main}. They motivated the analytic question answered
by the theorem: whether the centered primitive of the global
zero-counting discrepancy has sufficient unconditional control to
force the interior Li response below the natural \(H^{-2}\)
susceptibility scale.

The answer is affirmative by the classical estimate
\[
S_1(T)=O(\log T).
\]

\section{Horizontal blindness and scope}

Theorem~\ref{thm:main} is not an RH criterion.

To see the structural obstruction, consider two abstract zero
multisets having the same ordinates and multiplicities but different
real parts. Their scalar counting functions \(N(T)\) agree identically.
Therefore their discrepancies \(E(T)\), window means, and centered
primitives also agree.

The map
\[
\{\rho\}
\longmapsto
N(T)
\]
forgets the horizontal zero coordinate. Since RH is a statement about
\[
\operatorname{Re}\rho,
\]
an argument using only the scalar counting discrepancy cannot
distinguish all critical-line configurations from
ordinate-identical off-critical configurations.

Accordingly, the result proved here is a response-asymptotic theorem
for the scalar zero-counting discrepancy.

\section{Relation to classical Li and Weil machinery}

The kernel
\[
1-\left(1-\frac1s\right)^n
\]
is classical in Li's criterion. Bombieri and Lagarias
\cite{BombieriLagarias1999} derive an arithmetic formula for Li's
coefficients through the Guinand--Weil explicit formula and relate
Li positivity to Weil's criterion.

Likewise, the existence of an RH-equivalent positive quadratic
functional is classical Weil theory. Bombieri's study of Weil's
quadratic functional includes Fourier-side quadratic forms and finite
matrix truncations \cite{Bombieri2000}.

We do not claim a new Li criterion or a new Weil positivity framework.

The present paper instead studies the zero-coordinate Hessian
\[
2\operatorname{Re}
f_n''
\left(
\frac12+i\gamma
\right)
\]
in a matched-scale regime and combines its second-derivative profile
with a centered primitive of the classical zero-counting discrepancy.

No subsuming source has been identified in the current
uploaded-corpus audit for the exact combined response theorem stated
here. This sentence is not a novelty claim. Originality remains under active prior-art review.

\section{Conclusion}

We introduced the matched-scale Li susceptibility
\[
\LiSus_n(\gamma)
=
2\operatorname{Re}
f_n''
\left(
\frac12+i\gamma
\right)
\]
and derived compact-uniform profiles for the susceptibility and its
first two derivatives.

For the global zero-counting discrepancy
\[
E=N-M,
\]
we centered \(E\) on a fixed relative window and transferred its
response through the exact identity
\[
\Rint
=
-\Ebar
\left(
\LiSus_n(B)-\LiSus_n(H)
\right)
+
\int_H^B
F(t)\LiSus_n''(t)\,dt.
\]

The classical unconditional control
\[
S_1(T)=O(\log T)
\]
then yields
\[
\Rint
=
O_{K,\eta}
\left(
\frac{\log H}{H^3}
\right).
\]
Thus the interior discrepancy response is lower order than the
natural \(H^{-2}\) matched-scale susceptibility scale.

Under RH, the classical sharper first-primitive estimate supplies an
additional factor
\[
(\log\log H)^{-2}.
\]

The result is compatible with RH failure and should be interpreted as
an asymptotic transfer theorem rather than an RH criterion.

\begin{thebibliography}{99}

\bibitem{BombieriLagarias1999}
E. Bombieri and J. C. Lagarias,
\emph{Complements to Li's criterion for the Riemann hypothesis},
Journal of Number Theory \textbf{77} (1999), 274--287.

\bibitem{Bombieri2000}
E. Bombieri,
\emph{Remarks on Weil's quadratic functional in the theory of prime numbers, I},
Rendiconti Lincei. Matematica e Applicazioni,
Series 9, \textbf{11} (2000), 183--233.

\bibitem{WakasaS1}
T. Wakasa,
\emph{The explicit upper bound of the multiple integral of \(S(t)\)
on the Riemann Hypothesis}.
Used here as a source for the stated classical \(S_1\)-bound context
and references to the underlying classical results.

\end{thebibliography}

\end{document}
